\documentclass{article}
\usepackage[utf8]{inputenc}
\usepackage{booktabs}
\usepackage{siunitx}
\usepackage{caption}
\usepackage{threeparttable}
\usepackage{amsmath}

\sisetup{
    separate-uncertainty=true,
    table-align-uncertainty=true,
    table-format=1.2,
    tight-spacing=true,
    input-symbols = { ( ) },
    table-align-text-after=false
}

\begin{document}

\begin{table}[htbp]
    \centering
    \begin{threeparttable}
    \caption{Covariate Balance and Sample Characteristics}
    \label{tab:balance_table}
    \begin{tabular}{l *{4}{S}}
        \toprule
        & {Control} & {Treatment} & {Difference} & {$p$-value} \\
        Variable & {(N=450)} & {(N=442)} & {(1)-(2)} & {(t-test)} \\
        \midrule
        Age (years)      & 34.50  & 34.42  & 0.08  & 0.845 \\
                         & (8.20) & (8.15) &       &       \\
        Education (years) & 12.45  & 12.60  & -0.15 & 0.450 \\
                         & (2.40) & (2.35) &       &       \\
        Household Income & 45.20  & 46.10  & -0.90 & 0.312 \\
        (\$1,000s)       & (15.4) & (16.2) &       &       \\
        Female (\%)      & 0.52   & 0.51   & 0.01  & 0.910 \\
                         & (0.50) & (0.50) &       &       \\
        Urban (\%)       & 0.78   & 0.80   & -0.02 & 0.785 \\
                         & (0.41) & (0.40) &       &       \\
        \midrule
        \textbf{Joint orthogonality F-test (p-value)} & & & & 0.645 \\
        \bottomrule
    \end{tabular}
    \begin{tablenotes}
        \small
        \item \textit{Notes:} This table reports the mean values of baseline characteristics for the control and treatment groups. Standard deviations are in parentheses. The $p$-value in the last column is for a t-test of equality of means. The bottom row reports the $p$-value from a joint F-test of all covariates.
    \end{tablenotes}
    \end{threeparttable}
\end{table}

\end{document}
